\begin{explanationbox}
\textbf{一、先用一句话理解}

\textbf{圆的半径垂直于切线; 椭圆是把圆沿两个坐标方向按不同比例拉伸后的图形.}
圆上的斜率积是 $-1$; 横向拉得比纵向更长, 两条斜率都会缩小为原来的 $b/a$ 倍,
因此它们的乘积就变成 $-(b/a)^2$.

\textbf{二、为什么圆是 $-1$, 椭圆却不是?}

\begin{enumerate}
  \item 先看单位圆 $u^2+v^2=1$. 在非轴顶点 $Q(u_0,v_0)$ 处,
  半径 $OQ$ 与切线垂直, 所以两条斜率之积为 $-1$.
  \item 把圆上的点 $(u,v)$ 变成椭圆上的点 $(x,y)=(au,bv)$.
  这相当于横向放大 $a$ 倍, 纵向放大 $b$ 倍.
  \item 原来一条直线的斜率为 $m=\Delta v/\Delta u$; 变形后斜率为
  \[
  k=\frac{\Delta y}{\Delta x}
   =\frac{b\Delta v}{a\Delta u}=\frac ba m.
  \]
  \item 半径方向与切线方向\textbf{都经过同一种拉伸}, 因而两个斜率各乘一次 $b/a$:
  \[
  k_lk_{OP}=\left(\frac ba\right)^2(-1)=-\frac{b^2}{a^2}.
  \]
\end{enumerate}

\textbf{记忆图像}: \textbf{圆: $-1$; 椭圆: $-\text{短轴平方}/\text{长轴平方}$.}
这里假设长轴沿 $x$ 轴; 若长轴沿 $y$ 轴, 要按实际方程重新辨认系数.

\textbf{三、再从切线的法向量理解}

椭圆在 $P$ 处的切线满足
\[
\frac{x_0}{a^2}x+\frac{y_0}{b^2}y=1.
\]
一般直线 $Ax+By=C$ 的法向量为 $(A,B)$, 因为沿直线任意移动
$(\Delta x,\Delta y)$ 都满足 $A\Delta x+B\Delta y=0$.
因此这条切线的法向量可以取
\[
\boldsymbol n=\left(\frac{x_0}{a^2},\frac{y_0}{b^2}\right).
\]
而由中心指向切点的向量是 $\overrightarrow{OP}=(x_0,y_0)$.
\textbf{关键区别}: 在圆上两者方向相同; 在非圆椭圆的一般位置,
由于横纵坐标除以了不同的平方, $\boldsymbol n$ 通常不与 $OP$ 平行.
所以椭圆的切线\textbf{通常并不垂直}于 $OP$.
只有在四个轴顶点, 二者才垂直; 此时至少一条直线的斜率不存在.

\textbf{四、带数字看变化: $a=2$, $b=1$}

椭圆为 $x^2/4+y^2=1$. 取第一象限点
\[
P=\left(\sqrt2,\frac1{\sqrt2}\right).
\]
中心连线的斜率为
\[
k_{OP}=\frac{1/\sqrt2}{\sqrt2}=\frac12.
\]
切线斜率为
\[
k_l=-\frac{1\cdot\sqrt2}{4(1/\sqrt2)}=-\frac12.
\]
故 $k_lk_{OP}=-1/4$, 并非 $-1$.
两条直线恰好一上一下倾斜, 但\textbf{不是互相垂直}.

\textbf{五、定值到底说明什么?}

\begin{itemize}
  \item \textbf{点在变, 乘积不变}: $P$ 沿椭圆移动时, $k_l$ 与 $k_{OP}$ 各自变化,
  但只要两个斜率都存在, 乘积始终为 $-b^2/a^2$.
  \item \textbf{一正一负}: 因为乘积为负, 两斜率异号. 注意这不等于两直线垂直.
  \item \textbf{只与轴比有关}: 把 $a,b$ 同时放大同一倍, 斜率积不变.
  \item \textbf{与离心率有关}: $e^2=1-b^2/a^2$, 所以斜率积也可写成 $e^2-1$.
  \item \textbf{与中点弦有关}: 一组平行弦的中点沿一条过中心的直线排列;
  该中点连线的斜率与弦斜率也满足同样的乘积定值.
\end{itemize}

\textbf{六、做题时的识别信号}

看到“椭圆切线”“切点与中心的连线”“斜率积”或“平行弦中点”,
可优先想到 $k_lk_{OP}=-b^2/a^2$.
如果题目只求一个斜率, 常可直接使用定值;
如果还要求切线方程, 必须继续求切点坐标或切线的截距, 不能只写斜率.
\end{explanationbox}
